\documentclass[11pt]{article}
\usepackage[margin=1in]{geometry}
\usepackage{amsmath,amssymb}
\usepackage{booktabs}
\usepackage{hyperref}
\usepackage{enumitem}
\usepackage{xcolor}

\hypersetup{
  colorlinks=true,
  linkcolor=blue,
  citecolor=blue,
  urlcolor=blue
}

\title{A source-frozen Friedmann-readout audit protocol for late-time cosmology}
\author{Jozsef Olcsak}
\date{Draft protocol scaffold, 2026-06-08}

\begin{document}
\maketitle

\begin{abstract}
We define a reproducible audit protocol for a Tau-Core-inspired
Friedmann-readout residual hypothesis. The external motivation is the
Koushiappas cosmological uncertainty proposal, in which non-sharp
scale-factor and expansion-rate variables modify Friedmann closure and can
produce late-time acceleration without adding new particles or fields. The
present paper does not validate that model and does not claim that Tau Core
explains dark energy. Instead, it translates the idea into a source-frozen
cosmological audit: independent void-maturity and cosmic-web-anisotropy
source blocks define activation functions before distance, growth, lensing,
and CMB-safety observables are scored. The target signature is not merely
\(w(z)\ne -1\), but a Lambda-like distance readout paired with
web/environment-dependent growth or lensing residuals. The first release is a
protocol scaffold with claim boundaries, source blocks, null controls, and
reproducibility checks; no measurement result is claimed.
\end{abstract}

\section{Scope and non-claims}

This manuscript opens a cosmology-facing audit lane for Tau Core. It should be
read as a protocol scaffold, not as a discovery paper. It is downstream of the
finite-memory cosmology diagnostic in Paper 5~\cite{Paper5FiniteMemory} and
inherits leakage-prevention discipline from Papers 8 and 9
~\cite{Paper8ForwardReadout,Paper9ProjectionReadout}.

The allowed claim is:

\begin{quote}
A source-frozen cosmological readout audit can be formulated in which void and
web source proxies are frozen before testing distance, growth, lensing, and CMB
safety observables.
\end{quote}

The disallowed claims are:

\begin{itemize}[leftmargin=1.2em]
\item that Tau Core explains cosmic acceleration;
\item that the Koushiappas uncertainty-Friedmann model has been validated;
\item that \(\Lambda\)CDM is replaced;
\item that CMB maps directly show a Tau Core effect;
\item that a perturbation-level covariant Tau Core cosmology has been derived.
\end{itemize}

\section{External motivation}

Koushiappas~\cite{Koushiappas2026} proposes that the size of the universe and
its expansion rate cannot be simultaneously specified with arbitrary precision.
In that model a deformed commutation relation for the scale factor modifies the
Friedmann equation through a geometric correction. For a positive exponent the
model can produce late-time dark-energy-like behavior, while a sufficiently
negative exponent can produce a nonsingular bounce.

Tau Core does not import this as proof. It uses the paper as a comparison route:

\[
\begin{aligned}
&\hbox{non-sharp scale/expansion readout}\\
&\quad\longrightarrow\hbox{ modified Friedmann closure}\\
&\quad\longrightarrow\hbox{ late-time acceleration-like reconstruction}.
\end{aligned}
\]

\section{Friedmann-readout residual shell}

For a flat FLRW model without an explicit cosmological constant, define the
standard closure
\begin{equation}
C_F[a,H,\rho] = H^2 - \frac{8\pi G}{3}\rho.
\end{equation}

Tau Core introduces a bridge-level residual
\begin{equation}
C_{\tau,\mathrm{cosmo}}[a,H,\mathcal M_\tau],
\end{equation}
where \(a(t)\) and \(H(t)\) are four-dimensional cosmological readouts and
\(\mathcal M_\tau\) denotes a deeper cosmic morphology/projection/readout
structure. The claim-safe shell is
\begin{equation}
C_F[a,H,\rho]+C_{\tau,\mathrm{cosmo}}[a,H,\mathcal M_\tau]=0.
\end{equation}

Equivalently,
\begin{equation}
H^2=\frac{8\pi G}{3}\rho-C_{\tau,\mathrm{cosmo}}.
\end{equation}

If
\begin{equation}
-C_{\tau,\mathrm{cosmo}}\simeq \frac{\Lambda_{\mathrm{eff}}}{3},
\end{equation}
then an observer reconstructs a Lambda-like term. This is not a statement that
\(\Lambda_{\mathrm{eff}}\) is fundamental vacuum energy.

\section{Late-time activation}

A minimal late-activation shell is
\begin{equation}
C_{\tau,\mathrm{cosmo}}(z)
=
-\frac{\Lambda_{\tau,0}}{3}W_\tau(z),
\end{equation}
so that
\begin{equation}
H^2(z)
=
\frac{8\pi G}{3}\rho(z)
+\frac{\Lambda_{\tau,0}}{3}W_\tau(z).
\end{equation}

The safety behavior is
\[
W_\tau(z)\to 0 \quad (z\gg 1), \qquad
W_\tau(z)\to 1 \quad (z\to 0).
\]

One source-side activation shell is
\begin{equation}
W_\tau(a)=\frac{a^q}{a^q+a_\tau^q}.
\end{equation}

If a four-dimensional observer reconstructs the residual as an effective dark
energy component, then
\begin{equation}
w_{\mathrm{eff}}(z)
=
-1+\frac{1+z}{3}\frac{d\ln W_\tau}{dz}.
\end{equation}
This reconstructed equation of state is an observer-side diagnostic, not a
literal fluid claim.

\section{Void/web split target}

The audit should not stop at fitting \(H(z)\). The Tau-specific target is
\begin{equation}
\hbox{Lambda-like distance readout}
+
\hbox{environment-dependent growth/lensing residual}.
\end{equation}

The minimal distance shell is
\begin{equation}
E_D^2(a)
=
\Omega_m a^{-3}
+\Omega_r a^{-4}
+\Omega_k a^{-2}
+\Omega_\tau W_V(a),
\end{equation}
where \(W_V\) must be frozen from a void-maturity source block before distance
scoring.

The minimal growth shell is
\begin{equation}
D''
+
\left(2+\frac{d\ln H_D}{d\ln a}\right)D'
-
\frac{3}{2}\Omega_m(a)
\left[1+\mu_\tau(a,k,\mathrm{env})\right]D
=0,
\end{equation}
with
\begin{equation}
\mu_\tau(a,k,\mathrm{env})
=
\mu_0 W_W(a)Q_k(k)Q_{\mathrm{env}}(\mathrm{env}).
\end{equation}

The corresponding lensing shell is
\begin{equation}
\Sigma_{\mathrm{lens}}(a,k,\mathrm{env})
=
1+\Sigma_0 W_W(a)Q_k(k)Q_{\mathrm{env}}(\mathrm{env}).
\end{equation}

Here \(W_W\) must be frozen from a cosmic-web-anisotropy source block before
growth or lensing residuals are scored.

\section{Data blocks}

The intended public data blocks are summarized in
Table~\ref{tab:data-blocks}. The table is a work order, not a completed
measurement.

\begin{table}[h]
\centering
\small
\begin{tabular}{lll}
\toprule
Block & Candidate public sources & Role \\
\midrule
Distance & DESI BAO, SNe, \(H(z)\), CMB priors & test \(E_D(z)\) \\
Void maturity & void catalogues & freeze \(W_V(a)\) \\
Web anisotropy & cosmic-web classifiers & freeze \(W_W(a)\), \(Q_{\rm env}\) \\
Growth & RSD \(f\sigma_8\), clustering & test \(\mu_\tau\) \\
Lensing & weak lensing, CMB lensing, cluster lensing & test \(\Sigma_{\rm lens}\) \\
CMB safety & Planck/LAMBDA products & early safety and lensing controls \\
\bottomrule
\end{tabular}
\caption{Planned source blocks for the Friedmann-readout audit. DESI data
products~\cite{DESIDR2BAO2025,DESIDR2Products2025}, Planck products
~\cite{Planck2018Cosmology,PlanckLegacyArchive}, and NASA LAMBDA
~\cite{NASALAMBDA}, the public Cosmic Chronometers data page
~\cite{CosmicChronometersData}, SDSS cosmic-web and eBOSS/LSS products
~\cite{SDSSCosmicWebVAC,SDSSDR16LSS}, and KiDS-1000 lensing products
~\cite{KiDS1000Lensing} are examples of public source routes.}
\label{tab:data-blocks}
\end{table}

The repository also contains a source-freeze manifest:
\[
\texttt{data/source\_manifest/paper10\_source\_freeze\_manifest.csv}.
\]
It records candidate public source families, whether the source has merely been
selected or remains blocked, the planned local artifact path, and the endpoint
fields forbidden during source selection. In the current manifest, the first
source-selection audit has identified candidate public source routes for all
required blocks, but no source row is yet \texttt{FETCHED\_FROZEN}. This is the
key distinction: source-family selection is not endpoint readiness. These
blocks cannot be used for endpoint scoring until the public source product,
version, transformation rule, and checksum policy are frozen.

\section{Leakage prevention}

The audit follows a source-freeze discipline:

\begin{enumerate}[leftmargin=1.5em]
\item Freeze void maturity \(W_V\) before distance residual scoring.
\item Freeze web anisotropy \(W_W\) and \(Q_{\rm env}\) before growth/lensing scoring.
\item Keep CMB primary maps as early-safety and consistency controls.
\item Run shuffled environment labels and wrong-proxy nulls.
\item Do not use target residuals to choose activation functions.
\end{enumerate}

This is the cosmological analogue of the matched-vs-wrong and non-overlap
ledger discipline used in Papers 8 and 9.

In the current scaffold, the validation script
\(\texttt{scripts/validate\_source\_freeze.py}\) checks that the manifest covers
the required blocks and explicitly forbids endpoint leakage fields such as
distance residuals, growth residuals, lensing residuals, fit-improvement flags,
and post-hoc environment labels. Passing this script is not a data result; it is
a precondition for a later real-data pilot.

The repository further includes a compact provenance cache of the official
source landing pages. The acquisition script
\(\texttt{scripts/acquire\_source\_pages.py}\) records the landing-page payloads,
byte counts, and SHA-256 hashes, while
\(\texttt{scripts/validate\_cached\_source\_pages.py}\) verifies the committed
cache during reproduction. These cached pages are not raw survey products and
are not used to compute cosmological observables; they only make the source
routes auditable without requiring live network access during reproduction.

A second preflight artifact,
\(\texttt{data/source\_manifest/paper10\_product\_freeze\_plan.csv}\),
declares the product role, planned local path, transformation-rule identifier,
covariance policy, mask or selection policy, and forbidden endpoint fields for
each source family. In this scaffold every product row has
\(\texttt{endpoint\_allowed=false}\). Thus the protocol now specifies the next
data-acquisition step, but still does not contain endpoint data or a
measurement result.

\section{CMB role}

CMB maps are relevant, but not as the first standalone discovery target. The
primary CMB mostly constrains early-universe safety. More direct late-time
tests can come from CMB lensing, ISW-style cross-correlations, and
foreground-safe large-scale controls. Therefore the CMB role in this protocol is
\[
\hbox{early safety + lensing/growth cross-check},
\]
not direct proof of a Tau Core residual.

\section{Pass and fail gates}

A pass candidate requires:

\begin{itemize}[leftmargin=1.2em]
\item independent void maturity predicting distance activation;
\item independent web anisotropy predicting growth/lensing residuals;
\item degradation under shuffled environment labels;
\item survival under survey mask and covariance controls;
\item early-universe safety;
\item comparison against smooth dark energy and generic modified-gravity controls.
\end{itemize}

A fail result occurs if:

\begin{itemize}[leftmargin=1.2em]
\item \(W_V\) or \(W_W\) is reverse-engineered from target residuals;
\item no environment dependence appears;
\item smooth dark energy explains the data equally well;
\item the signal disappears under realistic mocks;
\item early-universe constraints fail.
\end{itemize}

\section{Conclusion}

The present repository defines a source-frozen Friedmann-readout audit protocol.
It does not report a measurement result. Its value is to make the next
cosmology test precise: a Tau Core-inspired late-time residual should be judged
not by \(w(z)\ne -1\) alone, but by the joint pattern
\[
\begin{aligned}
&\hbox{Lambda-like distance behavior}\\
&\quad+\hbox{ web/environment-dependent growth or lensing residual}\\
&\quad+\hbox{ CMB early-safety}.
\end{aligned}
\]

The next paper-grade step is a public-data pilot analysis with DESI/SNe/\(H(z)\)
distance data, void and web source proxies, growth/lensing blocks, CMB safety
checks, and predeclared null controls.

\bibliographystyle{plain}
\bibliography{refs}

\end{document}
